% =====================================================================
\chapter{Soluciones}
\label{ap:soluciones}
% =====================================================================

Las referencias a los apuntes siguen el formato «M21, §3.4»: módulo, y número de tema y sección tal como aparecen en el PDF de cada módulo. En los pasos intermedios se conservan al menos 4 cifras significativas y los resultados se dan con 3, como piden las instrucciones; cuando existe, se da también el valor exacto.

% ---------------------------------------------------------------------
\section{Parte A: clave de respuestas}

\begin{center}
\small
\renewcommand{\arraystretch}{1.12}
\begin{tabularx}{\textwidth}{@{}c c >{\raggedright\arraybackslash}X >{\raggedright\arraybackslash}p{2.5cm}@{}}
\toprule
Pregunta & Resp. & Justificación & Apuntes \\
\midrule
1 & (b) & La pérdida enmascarada solo puntúa la respuesta y su token de fin, imprescindible para que el modelo aprenda a terminar el turno. & M20, §2.2 y §2.4 \\
2 & (c) & $r(d+k)=8\cdot(4096+1024)=40\,960$. & M20, §5.1 \\
3 & (d) & Con $B=0$, $BA=0$; $\partial\mathcal{L}/\partial B=s\,G(Ax)\T\ne 0$ y $A$ recibe gradiente desde el segundo paso. & M20, §5.3 \\
4 & (a) & $\theta^{\ast}=(a+\lambda\theta_0)/(1+\lambda)=4/4=1$. & M20, §3.2 \\
5 & (b) & $b+32/B=4+32/32=5$ bits. & M20, §6.1 \\
6 & (d) & $G_0=R_1+\gamma R_2+\gamma^{2}R_3=1+0+0{,}25\cdot 2=1{,}5$. & M21, §1.2 \\
7 & (b) & La puntuación tiene media nula, de modo que $\E[\nabla\ln\pi\cdot b(s)]=0$. & M21, §2.5 \\
8 & (c) & Con $A<0$, el término es $\max(\rho,1-\varepsilon)A=0{,}8\cdot(-1)=-0{,}8$. & M21, §3.2 \\
9 & (d) & $\widehat{r_a-r_b}=\ln\bigl(n_a/(n-n_a)\bigr)=\ln(8/2)=\ln 4$. & M21, §4.2 \\
10 & (d) & Los dos log-cocientes son nulos y $-\ln\sigma(0)=\ln 2$. & M21, §5.3 \\
11 & (a) & Todos los numeradores $r_i-\bar r$ son nulos; con $\varepsilon_0>0$ las ventajas son~0. & M21, §6.3 \\
12 & (b) & $\idf_t=\ln(1000/10)=\ln 100=4{,}605$; el valor~2 corresponde al logaritmo decimal. & M22, §2.2 \\
13 & (c) & $g(f,\ell)\to k_1+1=2{,}5$: la frecuencia se satura. & M22, §2.6 \\
14 & (a) & $\AP=\frac14\bigl(\frac12+\frac24+\frac35\bigr)=1{,}6/4=0{,}4$; (c) divide entre 3, no entre $\abs{\Rel}$. & M22, §3.2 \\
15 & (a) & Si el modelo ha visto el examen, la puntuación mide memoria y no capacidad. & M22, §3.4 \\
16 & (c) & $0{,}98^{20}=0{,}668$; (b) es la aproximación lineal $1-20\cdot 0{,}02$. & M22, §6.4 \\
17 & (d) & Las características dispersas casi nunca coinciden y la interferencia se paga rara vez. & M23, §3.1--3.3 \\
18 & (a) & $a-\lambda/2=0{,}3-0{,}4<0$, luego $f^{\ast}=0$: la $L^{1}$ anula las activaciones pequeñas. & M23, §4.3 \\
19 & (b) & Dividir por $T>0$ conserva el orden de los logits y, por tanto, el argmax. & M23, §6.3 \\
20 & (c) & Regla del 3: $3/600=0{,}5\,\%$ (la cota exacta es $0{,}498\,\%$). & M23, §6.6 \\
\bottomrule
\end{tabularx}
\end{center}

% ---------------------------------------------------------------------
\section{Parte B: soluciones de los ejercicios}

% =====================================================================  B.1
\begin{solucion}{eje:p1}
\apartado{(a)}
\datos instrucciones $[12,7,3,9]$ y $[6,2]$; respuestas $[4,0]$ y $[8,5,0]$; relleno~$1$.
\planteamiento se concatena, se rellena a la longitud máxima y se desplaza una posición: la etiqueta de la posición $t$ es el token $t+1$. La máscara vale~1 cuando esa etiqueta es un token de la respuesta (M20, §2.2 y~§2.3).
\pasos
\begin{enumerate}[label=\arabic*.]
  \item Secuencias: $[12,7,3,9,4,0]$ (6 tokens) y $[6,2,8,5,0]$ (5 tokens), que se rellena a $[6,2,8,5,0,1]$.
  \item Ejemplo~1: \texttt{input\_ids} $=[12,7,3,9,4]$, \texttt{labels} $=[7,3,9,4,0]$. Las etiquetas $4$ y $0$ son de la respuesta: máscara $[0,0,0,1,1]$.
  \item Ejemplo~2: \texttt{input\_ids} $=[6,2,8,5,0]$, \texttt{labels} $=[2,8,5,0,1]$. Las etiquetas $8$, $5$ y $0$ son de la respuesta y la última es relleno: máscara $[0,1,1,1,0]$.
\end{enumerate}
\resultado \textbf{puntúan 2 posiciones en el ejemplo~1 y 3 en el ejemplo~2}, tantas como tokens tiene cada respuesta.
\comprobacion la última etiqueta de la instrucción ($9$ en el ejemplo~1) no puntúa, y la primera posición que puntúa es la que tiene como entrada el último token de la instrucción ($9\to 4$).
\criterio 1 punto por las secuencias desplazadas; 1,5 por las máscaras; 0,5 por el recuento.

\apartado{(b)}
\datos pérdidas $(2{,}4;\ 1{,}9;\ 0{,}7;\ 0{,}8;\ 0{,}2)$ y $(1{,}6;\ 0{,}9;\ 0{,}6;\ 0{,}3;\ 3{,}0)$; máscaras de (a); $B=2$; $c=5$.
\planteamiento sumas enmascaradas $S_i=\sum_tm^{(i)}_t\ell^{(i)}_t$ y los cuatro cocientes del enunciado; la acumulación es exacta si la pérdida del lote es una suma de términos por secuencia con denominadores fijos (M20, §2.3, proposición de acumulación de gradientes).
\pasos
\begin{enumerate}[label=\arabic*.]
  \item $S_1=0{,}8+0{,}2=1{,}0$ (2 posiciones) y $S_2=0{,}9+0{,}6+0{,}3=1{,}8$ (3 posiciones).
  \item (i) $J_{\mathrm{tok}}=(1{,}0+1{,}8)/(2+3)=2{,}8/5=0{,}56$.
  \item (ii) $J_{\mathrm{mm}}=\frac12\bigl(1{,}0/2+1{,}8/3\bigr)=\frac12(0{,}5+0{,}6)=0{,}55$.
  \item (iii) $J_{\mathrm{c}}=(1{,}0+1{,}8)/(2\cdot 5)=0{,}28$; es la normalización $J_{\mathrm{sec}}$ de los apuntes.
  \item (iv) Sin máscara, excluyendo el relleno: $(2{,}4+1{,}9+0{,}7+0{,}8+0{,}2)+(1{,}6+0{,}9+0{,}6+0{,}3)=6{,}0+3{,}4=9{,}4$ sobre 9 posiciones, $9{,}4/9=1{,}044$.
  \item Acumulación con (i): cada microlote divide entre sus propios tokens, $1{,}0/2=0{,}5$ y $1{,}8/3=0{,}6$; la suma de $J_g/2$ es $0{,}55\neq 0{,}56$.
  \item Acumulación con (iii): cada microlote calcula $S_g/(1\cdot 5)$ y retropropaga la mitad, $0{,}1+0{,}18=0{,}28=J_{\mathrm{c}}$.
\end{enumerate}
\resultado \textbf{$J_{\mathrm{tok}}=0{,}56$; $J_{\mathrm{mm}}=0{,}55$; $J_{\mathrm{c}}=0{,}28$; sin máscara, $1{,}04$. La acumulación reproduce $J_{\mathrm{c}}$, no $J_{\mathrm{tok}}$}: la suma de cocientes no es el cociente de las sumas cuando los microlotes tienen distinto número de tokens que puntúan; con un denominador constante, la pérdida es lineal en las secuencias y el gradiente también.
\comprobacion la pérdida sin máscara es mucho mayor porque la instrucción, que el modelo no tiene por qué predecir, domina la suma ($6{,}6$ de $9{,}4$).
\criterio 0,5 puntos por cada pérdida (i)--(iv); 1 punto por la acumulación y su justificación.

\apartado{(c)}
\datos longitudes $11$, $6$, $14$ y $8$; un token de fin por ejemplo; $m=8$; relleno hasta 16.
\planteamiento con empaquetado se obtienen $\lfloor T/m\rfloor$ secuencias de los $T$ tokens totales; con relleno, la fracción útil es $T/(4\cdot 16)$ (M20, §2.3).
\pasos
\begin{enumerate}[label=\arabic*.]
  \item Con el token de fin: $12+7+15+9=43$ tokens.
  \item $\lfloor 43/8\rfloor=5$ secuencias, con $40$ tokens; se pierden $43-40=3$ tokens, el $7{,}0\,\%$.
  \item Con relleno: $4\cdot 16=64$ posiciones, de las que son útiles $43$: $43/64=67{,}2\,\%$.
\end{enumerate}
\resultado \textbf{5 secuencias y 3 tokens perdidos; con relleno solo el $67{,}2\,\%$ de las posiciones es útil.} Inconveniente: una secuencia empaquetada mezcla el final de un ejemplo con el principio de otro, y la máscara causal permite que un ejemplo atienda al anterior; además, los ejemplos quedan cortados.
\comprobacion los límites entre ejemplos están en las posiciones $12$, $19$, $34$ y $43$; la segunda secuencia ($8$ a $15$) contiene el final del primer ejemplo y el principio del segundo.
\criterio 1 punto por el empaquetado; 0,5 por el relleno; 0,5 por el inconveniente.

\apartado{(d)}
\datos $x_{1:T}=(u_{1:n},y_{1:m})$; $m_t=\mathbb{1}\{t>n\}$.
\planteamiento regla de la cadena condicionada a $u$ y gradiente de $\LSE$ (M20, §2.2).
\pasos
\begin{enumerate}[label=\arabic*.]
  \item $\sum_tm_t=m$ y $\sum_tm_t\ln p_\theta(x_t\mid x_{1:t-1})=\sum_{j=1}^{m}\ln p_\theta(y_j\mid u,y_{1:j-1})=\ln p_\theta(y\mid u)$, por la regla de la cadena.
  \item Dividiendo por $m$ y cambiando el signo, $\Lsft=-\frac1m\ln p_\theta(y\mid u)$: el SFT es máxima verosimilitud de la respuesta dada la instrucción.
  \item Solo el sumando $t$ depende de los logits $z_t$, y $\partial(-\ln p_t)/\partial z_t=p_t-e_{x_t}$; por tanto $\partial\Lsft/\partial z_t=\frac{m_t}{m}(p_t-e_{x_t})$, que es nulo si $t\le n$.
  \item La instrucción sí influye: las posiciones de la respuesta atienden a las de la instrucción, y el gradiente llega a sus estados ocultos y a los parámetros a través de la atención. Lo que se suprime es la señal que obligaría a \emph{predecir} los tokens de la instrucción.
\end{enumerate}
\resultado \textbf{$\Lsft=-\frac1m\ln p_\theta(y\mid u)$; gradiente nulo en los logits de la instrucción, pero no en sus representaciones.}
\comprobacion con $m_t=1$ para todo $t$ se recupera la pérdida del preentrenamiento, $-\frac1T\ln p_\theta(u,y)$.
\criterio 1 punto por la identidad; 0,5 por el gradiente; 0,5 por la discusión.
\repasar{M20, §2.2, §2.3 y §2.4.}
\end{solucion}

% =====================================================================  B.2
\begin{solucion}{eje:p2}
\apartado{(a)}
\datos $W_0=\begin{pmatrix}2&0\\0&1\end{pmatrix}$, $B=(1,2)\T$, $A=(2,-1)$, $r=1$, $\alpha=0{,}5$, $x=(1,3)\T$, $G=(1,0)\T$.
\planteamiento $s=\alpha/r$; $h=W_0x+s\,B(Ax)$; $W'=W_0+s\,BA$; $\partial\mathcal{L}/\partial B=s\,G(Ax)\T$ y $\partial\mathcal{L}/\partial A=s\,B\T G\,x\T$ (M20, §5.1 y~§5.3).
\pasos
\begin{enumerate}[label=\arabic*.]
  \item $s=0{,}5/1=0{,}5$; $Ax=2\cdot 1-1\cdot 3=-1$; $B(Ax)=(-1,-2)\T$.
  \item $h=(2,3)\T+0{,}5\cdot(-1,-2)\T=(1{,}5;\ 2)\T$.
  \item $BA=\begin{pmatrix}2&-1\\4&-2\end{pmatrix}$, de rango~1; $W'=\begin{pmatrix}2+1&-0{,}5\\0+2&1-1\end{pmatrix}=\begin{pmatrix}3&-0{,}5\\2&0\end{pmatrix}$.
  \item $W'x=(3-1{,}5;\ 2+0)\T=(1{,}5;\ 2)\T=h$.
  \item $\partial\mathcal{L}/\partial B=0{,}5\cdot(1,0)\T\cdot(-1)=(-0{,}5;\ 0)\T$.
  \item $B\T G=1\cdot 1+2\cdot 0=1$, luego $\partial\mathcal{L}/\partial A=0{,}5\cdot 1\cdot(1,3)=(0{,}5;\ 1{,}5)$.
  \item Con $B=0$, $B\T G=0$ y $\partial\mathcal{L}/\partial A=(0,0)$.
\end{enumerate}
\resultado \textbf{$h=W'x=(1{,}5;\ 2)$; $W'=\begin{pmatrix}3&-0{,}5\\2&0\end{pmatrix}$; $\partial\mathcal{L}/\partial B=(-0{,}5;\ 0)\T$, $\partial\mathcal{L}/\partial A=(0{,}5;\ 1{,}5)$; con $B=0$, $\partial\mathcal{L}/\partial A=0$.}
\comprobacion las formas son las de los parámetros: $\partial\mathcal{L}/\partial B$ es $2\times 1$ y $\partial\mathcal{L}/\partial A$ es $1\times 2$. La fusión no añade coste de inferencia.
\criterio 1 punto por $h$ y la fusión; 1,5 por los gradientes; 0,5 por el caso $B=0$.

\apartado{(b)}
\datos $\sigma=(4;\ 2;\ 1;\ 0{,}5)$.
\planteamiento por Eckart--Young, el menor error de rango $r$ es $\bigl(\sum_{i>r}\sigma_i^{2}\bigr)^{1/2}$ (M20, §5.5).
\pasos
\begin{enumerate}[label=\arabic*.]
  \item Energía total: $16+4+1+0{,}25=21{,}25$.
  \item $r=1$: error $\sqrt{4+1+0{,}25}=\sqrt{5{,}25}=2{,}291$; energía $16/21{,}25=75{,}3\,\%$.
  \item $r=2$: error $\sqrt{1{,}25}=1{,}118$; energía $20/21{,}25=94{,}1\,\%$.
\end{enumerate}
\resultado \textbf{errores $2{,}29$ y $1{,}12$; energías $75{,}3\,\%$ y $94{,}1\,\%$; el menor rango con al menos el $90\,\%$ es $r=2$.}
\comprobacion error$^{2}$ + energía conservada $=5{,}25+16=1{,}25+20=21{,}25$.

\apartado{(c)}
\datos 32 capas; 4 matrices $4096\times 4096$ por capa; $r=16$; $N=7\cdot 10^{9}$; $B_1=64$, $B_2=256$.
\planteamiento $r(d+k)$ por matriz; $b_{\mathrm{ef}}=4+8/B_1+32/(B_1B_2)$; memoria $(b_{\mathrm{ef}}/8)N+16N_{\mathrm{LoRA}}$ frente a $16N$ (M20, §5.1, §6.3 y~§6.4).
\pasos
\begin{enumerate}[label=\arabic*.]
  \item Por matriz: $16\cdot 8192=131\,072$. Total: $131\,072\cdot 4\cdot 32=16\,777\,216\approx 1{,}68\cdot 10^{7}$.
  \item Fracción: $2r/d=32/4096=0{,}781\,\%$ de las $2{,}15\cdot 10^{9}$ entradas de esas matrices.
  \item $b_{\mathrm{ef}}=4+0{,}125+0{,}00195=4{,}127$ bits.
  \item Base: $7\cdot 10^{9}\cdot 4{,}127/8=3{,}611\cdot 10^{9}$ B $=3{,}61$~GB. LoRA: $16\,777\,216\cdot 16=2{,}684\cdot 10^{8}$ B $=0{,}268$~GB. Total: $3{,}88$~GB.
  \item Ajuste completo: $16\cdot 7\cdot 10^{9}=112$~GB; cociente $112/3{,}88=28{,}9$.
\end{enumerate}
\resultado \textbf{$1{,}68\cdot 10^{7}$ parámetros ($0{,}781\,\%$); $b_{\mathrm{ef}}=4{,}127$ bits; QLoRA $\approx 3{,}88$~GB frente a $112$~GB, unas 29 veces menos.}
\comprobacion sin doble cuantización ($4{,}5$ bits) la base ocuparía $3{,}94$~GB; el ahorro de $0{,}373$ bits por parámetro son unos $0{,}33$~GB.
\criterio 1 punto por el recuento; 1 por $b_{\mathrm{ef}}$; 1 por la memoria y la comparación.

\apartado{(d)}
\datos $\theta_0=2$, $a=5$.
\planteamiento $\theta^{\ast}=(a+\lambda\theta_0)/(1+\lambda)$; olvido $(a-\theta_0)^{2}/(1+\lambda)^{2}$; error $\lambda^{2}(a-\theta_0)^{2}/(1+\lambda)^{2}$ (M20, §3.2).
\pasos
\begin{enumerate}[label=\arabic*.]
  \item $\lambda=2$: $\theta^{\ast}=(5+4)/3=3$; olvido $(3-2)^{2}=1$; error $(3-5)^{2}=4$.
  \item $9/(1+\lambda)^{2}\le 0{,}25\iff(1+\lambda)^{2}\ge 36\iff\lambda\ge 5$.
  \item $\lambda=5$: $\theta^{\ast}=15/6=2{,}5$; error $(2{,}5-5)^{2}=6{,}25$.
\end{enumerate}
\resultado \textbf{$\lambda=2$: $\theta^{\ast}=3$, olvido~1, error~4. Hace falta $\lambda\ge 5$, y el error en la tarea nueva sube a $6{,}25$.}
\comprobacion las distancias suman $\abs{a-\theta_0}=3$ en ambos casos: $1+2$ y $0{,}5+2{,}5$.
\repasar{M20, §3.2, §5.1, §5.3, §5.5, §6.1, §6.3 y §6.4.}
\end{solucion}

% =====================================================================  B.3
\begin{solucion}{eje:p3}
\apartado{(a)}
\datos MDP determinista de dos estados; $\gamma=0{,}5$; $\pi(S_1)=$ cambiar, $\pi(S_2)=$ quedarse.
\planteamiento ecuación de Bellman $V^{\pi}(s)=r(s,\pi(s))+\gamma V^{\pi}(s')$ y su versión iterativa $V_{k+1}=T^{\pi}V_k$ (M21, §1.3).
\pasos
\begin{enumerate}[label=\arabic*.]
  \item $V^{\pi}(S_2)=3+0{,}5\,V^{\pi}(S_2)\Rightarrow V^{\pi}(S_2)=6$; $V^{\pi}(S_1)=-2+0{,}5\cdot 6=1$.
  \item Iteración: $V_{k+1}(S_1)=-2+0{,}5\,V_k(S_2)$ y $V_{k+1}(S_2)=3+0{,}5\,V_k(S_2)$.
\end{enumerate}
\begin{center}
\small
\begin{tabular}{@{}ccccc@{}}
\toprule
$k$ & $V_k$ & error en $S_1$ & error en $S_2$ & máximo \\
\midrule
0 & $(0;\ 0)$ & $1$ & $6$ & $6$ \\
1 & $(-2;\ 3)$ & $3$ & $3$ & $3$ \\
2 & $(-0{,}5;\ 4{,}5)$ & $1{,}5$ & $1{,}5$ & $1{,}5$ \\
3 & $(0{,}25;\ 5{,}25)$ & $0{,}75$ & $0{,}75$ & $0{,}75$ \\
\bottomrule
\end{tabular}
\end{center}
\begin{enumerate}[label=\arabic*.,resume]
  \item El error en $S_1$ pasa de $1$ a $3$. No hay contradicción: la contracción afirma $\norm{T^{\pi}V-T^{\pi}V'}_\infty\le\gamma\norm{V-V'}_\infty$, una cota sobre el \emph{máximo}, que sí se divide por~2 en cada paso ($6\to 3\to 1{,}5\to 0{,}75$). En un estado concreto, $V_{k+1}(S_1)-V^{\pi}(S_1)=\gamma\bigl(V_k(S_2)-V^{\pi}(S_2)\bigr)$: el error de $S_1$ hereda el de $S_2$, que al principio era grande.
  \item $Q^{\pi}(S_1,\text{quedarse})=-1+0{,}5\,V^{\pi}(S_1)=-0{,}5$; $A^{\pi}(S_1,\text{quedarse})=-0{,}5-1=-1{,}5$.
  \item En $S_2$, $Q^{\pi}(S_2,\text{cambiar})=-1+0{,}5\cdot 1=-0{,}5<Q^{\pi}(S_2,\text{quedarse})=6$. La política voraz respecto de $Q^{\pi}$ coincide con $\pi$ en ambos estados; por el teorema de mejora de la política, $\pi$ es óptima (Sutton y Barto, §4.2~\cite{sb}; ampliación respecto de los apuntes).
\end{enumerate}
\resultado \textbf{$V^{\pi}=(1;\ 6)$; el error en $S_1$ sube de 1 a 3 mientras el máximo se divide por~2; $A^{\pi}(S_1,\text{quedarse})=-1{,}5$; $\pi$ es óptima.}
\comprobacion la iteración de valores óptima desde $V_0=0$ converge al mismo $V^{\ast}=(1;\ 6)$, y $\sum_a\pi(a\mid S_1)A^{\pi}(S_1,a)=A^{\pi}(S_1,\text{cambiar})=0$.
\criterio 1 punto por $V^{\pi}$; 1,5 por la tabla y la explicación de la contracción; 1,5 por la ventaja y la optimalidad.

\apartado{(b)}
\datos $\pi_\theta(1)=p=\sigma(\theta)$; $r(a)=\mathbb{1}\{a=1\}$; $b$ constante.
\planteamiento cada sumando $Z=(r(a)-b)\,g(a)$, con $g(a)=\frac{d}{d\theta}\ln\pi_\theta(a)$, toma dos valores; $\Var(\hat g_b)=\Var(Z)/n$ por independencia (M21, §2.5).
\pasos
\begin{enumerate}[label=\arabic*.]
  \item $g(1)=\frac{d}{d\theta}\ln\sigma(\theta)=1-p$ y $g(0)=\frac{d}{d\theta}\ln(1-\sigma(\theta))=-\sigma(\theta)(1-\sigma(\theta))/(1-\sigma(\theta))=-p$.
  \item $Z=(1-b)(1-p)$ con probabilidad $p$ y $Z=(-b)(-p)=bp$ con probabilidad $1-p$. Entonces $\E[Z]=p(1-p)(1-b)+(1-p)\,bp=p(1-p)$, que es $dJ/d\theta$ con $J=p$.
  \item $\E[Z^{2}]=p(1-p)^{2}(1-b)^{2}+(1-p)p^{2}b^{2}$, y
  \[
  \Var(Z)=p(1-p)\bigl[(1-p)(1-b)^{2}+pb^{2}-p(1-p)\bigr]=p(1-p)(1-p-b)^{2},
  \]
  porque el corchete vale $(1-p)-2b(1-p)+b^{2}-p+p^{2}=(1-p)^{2}-2b(1-p)+b^{2}$.
  \item $b=0$: $\Var=p(1-p)^{3}/n$. $b=p$: $\Var=p(1-p)(1-2p)^{2}/n$.
  \item El cociente es $(1-2p)^{2}/(1-p)^{2}<1\iff\abs{1-2p}<1-p$. Si $p\le\frac12$ se cumple siempre; si $p>\frac12$ equivale a $2p-1<1-p$, es decir, $p<\frac23$.
  \item $\Var(b)$ es una parábola en $b$ con mínimo en $b^{*}=1-p$, donde vale~0: $Z=p(1-p)$ para las dos acciones.
\end{enumerate}
\resultado \textbf{$\Var(\hat g_b)=p(1-p)(1-p-b)^{2}/n$; la línea base $b=p$ reduce la varianza solo si $p<2/3$ (la iguala en $p=2/3$ y la aumenta si $p>2/3$); la óptima es $b^{*}=1-p$, con varianza nula.}
\comprobacion con $p=0{,}25$ y $n=16$: $6{,}59\cdot 10^{-3}$ con $b=0$ y $2{,}93\cdot 10^{-3}$ con $b=p$, de acuerdo con $p<2/3$.
\criterio 1 punto por (i); 1,5 por (ii); 1 por (iii); 0,5 por (iv).

\apartado{(c)}
\datos $p=0{,}8$; $n=16$.
\planteamiento se sustituye en la fórmula de (b) y en $b^{*}=\E[rg^{2}]/\E[g^{2}]$ (M21, §2.5, proposición de la línea base de varianza mínima).
\pasos
\begin{enumerate}[label=\arabic*.]
  \item $b=0$: $0{,}8\cdot 0{,}2^{3}/16=0{,}0064/16=4{,}0\cdot 10^{-4}$.
  \item $b=0{,}8$: $0{,}8\cdot 0{,}2\cdot 0{,}6^{2}/16=0{,}0576/16=3{,}6\cdot 10^{-3}$, nueve veces más.
  \item $b^{*}=0{,}2$: varianza~0.
  \item $\E[rg^{2}]=p(1-p)^{2}=0{,}032$; $\E[g^{2}]=p(1-p)^{2}+(1-p)p^{2}=0{,}032+0{,}128=0{,}16$; $b^{*}=0{,}032/0{,}16=0{,}2=1-p$.
\end{enumerate}
\resultado \textbf{$4{,}0\cdot 10^{-4}$, $3{,}6\cdot 10^{-3}$ y $0$.} La recompensa media suele ser una buena línea base, pero no la óptima: aquí la acción rara ($a=0$) tiene una puntuación grande, $\abs{g(0)}=0{,}8$, y $b=p$ amplifica su contribución.
\comprobacion la esperanza no depende de $b$: $\E[\hat g_b]=0{,}8\cdot 0{,}2=0{,}16$ en los tres casos.
\repasar{M21, §1.3, §2.3 y §2.5.}
\end{solucion}

% =====================================================================  B.4
\begin{solucion}{eje:p4}
\apartado{(a)}
\datos $n=20$ comparaciones, $n_a=15$; razón $2{:}1$ entre $y_b$ e $y_c$.
\planteamiento estimador con dos objetos $\widehat{r_a-r_b}=\ln\bigl(n_a/(n-n_a)\bigr)$; las razones de probabilidades se multiplican (M21, §4.1 y~§4.2).
\pasos
\begin{enumerate}[label=\arabic*.]
  \item $\widehat{r_a-r_b}=\ln(15/5)=\ln 3=1{,}099$, y $\sigma(\ln 3)=3/4$.
  \item Pérdida media: $-\bigl(0{,}75\ln 0{,}75+0{,}25\ln 0{,}25\bigr)=0{,}2158+0{,}3466=0{,}5623$ nats.
  \item Razón $a{:}c=3\cdot 2=6$, luego $P(y_a\pref y_c)=6/7=0{,}857$.
\end{enumerate}
\resultado \textbf{$r_a-r_b=\ln 3\approx 1{,}10$; pérdida $0{,}562$ nats; $P(y_a\pref y_c)=6/7\approx 0{,}857$.}
\comprobacion la pérdida en el óptimo es la entropía binaria de la frecuencia observada, menor que $\ln 2=0{,}693$.

\apartado{(b)}
\datos $\piref=(0{,}4;\ 0{,}4;\ 0{,}2)$, $r=(1,0,2)$, $\beta\in\{1;\ 0{,}5\}$.
\planteamiento $\pi^{*}=\piref e^{r/\beta}/Z$; $\ln(\pi^{*}/\piref)=r/\beta-\ln Z$; el objetivo vale $\beta\ln Z$ (M21, §3.4).
\pasos
\begin{enumerate}[label=\arabic*.]
  \item $\beta=1$: pesos $(0{,}4e;\ 0{,}4;\ 0{,}2e^{2})=(1{,}0873;\ 0{,}4;\ 1{,}4778)$; $Z=2{,}9651$.
  \item $\pi^{*}=(0{,}3667;\ 0{,}1349;\ 0{,}4984)$; $\E_{\pi^{*}}[r]=0{,}3667+2\cdot 0{,}4984=1{,}3635$, frente a $0{,}8$ con $\piref$.
  \item $\ln(\pi^{*}/\piref)=(-0{,}0869;\ -1{,}0869;\ 0{,}9131)$; $\KL=-0{,}0319-0{,}1466+0{,}4551=0{,}2766$ nats.
  \item Objetivo: $1{,}3635-0{,}2766=1{,}0869=\ln 2{,}9651$.
  \item $\beta=0{,}5$: pesos $(2{,}9556;\ 0{,}4;\ 10{,}9196)$; $Z=14{,}275$; $\pi^{*}=(0{,}2070;\ 0{,}0280;\ 0{,}7649)$; $\E[r]=1{,}7369$; $\KL=0{,}8153$; objetivo $1{,}7369-0{,}5\cdot 0{,}8153=1{,}3293=0{,}5\ln 14{,}275$.
\end{enumerate}
\resultado \textbf{$\beta=1$: $\pi^{*}=(0{,}367;\ 0{,}135;\ 0{,}498)$, $Z=2{,}97$, objetivo $1{,}087=\ln Z$. $\beta=0{,}5$: $\pi^{*}=(0{,}207;\ 0{,}028;\ 0{,}765)$, objetivo $1{,}329=0{,}5\ln Z$.} Un $\beta$ menor permite alejarse más de la referencia (KL de $0{,}277$ a $0{,}815$) a cambio de más recompensa.
\comprobacion $\KL=\E_{\pi^{*}}[r]/\beta-\ln Z$: $1{,}3635-1{,}0869=0{,}2766$ y $1{,}7369/0{,}5-2{,}6585=0{,}8153$.
\criterio 2 puntos por el caso $\beta=1$ con la comprobación; 1 por el caso $\beta=0{,}5$ y la interpretación.

\apartado{(c)}
\datos $\pi^{*}$ en forma cerrada; $P(y_w\pref y_l\mid x)=\sigma\bigl(r(x,y_w)-r(x,y_l)\bigr)$.
\planteamiento despejar $r$, sustituir en Bradley--Terry, reemplazar $\pi^{*}$ por $\pi_\theta$ y maximizar la verosimilitud (M21, §5.1 y~§5.3).
\pasos
\begin{enumerate}[label=\arabic*.]
  \item Tomando logaritmos, $r(x,y)=\beta\ln\dfrac{\pi^{*}(y\mid x)}{\piref(y\mid x)}+\beta\ln Z(x)$.
  \item En $r(x,y_w)-r(x,y_l)$ el término $\beta\ln Z(x)$ aparece con los dos signos y se cancela, porque solo depende de $x$, común a las dos respuestas.
  \item Con $\pi_\theta$ en lugar de $\pi^{*}$ y la log-verosimilitud negativa de las preferencias:
  \[
  \LDPO(\theta)=-\E\Bigl[\ln\sigma\Bigl(\beta\ln\frac{\pi_\theta(y_w\mid x)}{\piref(y_w\mid x)}-\beta\ln\frac{\pi_\theta(y_l\mid x)}{\piref(y_l\mid x)}\Bigr)\Bigr].
  \]
  \item Por la regla de la cadena, $\ln\pi(x\oplus y)=\ln\pi(x)+\ln\pi(y\mid x)$; en la diferencia entre $x\oplus y_w$ y $x\oplus y_l$ se cancela $\ln\pi(x)$, para $\pi_\theta$ y para $\piref$.
\end{enumerate}
\resultado \textbf{la pérdida DPO anterior}: $Z(x)$ desaparece porque Bradley--Terry solo usa diferencias de recompensa para la misma instrucción. No hace falta un modelo de recompensa ni muestrear durante el entrenamiento.
\comprobacion con $\pi_\theta=\piref$ el argumento de $\sigma$ es~0 y la pérdida vale $\ln 2$, como en la cuestión~10.

\apartado{(d)}
\datos log-probabilidades $-20$, $-21{,}5$, $-24$ y $-23$; $\beta=0{,}1$.
\planteamiento $\hat r=\beta\ln(\pi_\theta/\piref)$; margen $h=\hat r_w-\hat r_l$; pérdida $-\ln\sigma(h)$; peso $\beta\,\sigma(-h)$ (M21, §5.3).
\pasos
\begin{enumerate}[label=\arabic*.]
  \item Log-cocientes: $-20+21{,}5=1{,}5$ y $-24+23=-1$; $\hat r_w=0{,}15$ y $\hat r_l=-0{,}10$.
  \item $h=0{,}25$; $\sigma(0{,}25)=0{,}5622$; pérdida $-\ln 0{,}5622=0{,}5759$.
  \item Peso: $0{,}1\cdot\sigma(-0{,}25)=0{,}1\cdot 0{,}4378=0{,}0438$.
\end{enumerate}
\resultado \textbf{$\hat r_w=0{,}15$, $\hat r_l=-0{,}10$, margen $0{,}25$, pérdida $0{,}576$ y peso $0{,}0438$.}
\comprobacion la pérdida es menor que $\ln 2$ porque el par ya está bien ordenado; el peso disminuiría si el margen creciera.

\apartado{(e)}
\datos $(1,0,0,1,0)$; probabilidad de acierto $0{,}4$.
\planteamiento $\widehat A_i=(r_i-\bar r)/\mathrm{desv}$; un grupo no aporta señal si todas las recompensas coinciden (M21, §6.3).
\pasos
\begin{enumerate}[label=\arabic*.]
  \item $\bar r=0{,}4$; desviación poblacional $\sqrt{0{,}24}=0{,}4899$; ventajas $0{,}6/0{,}4899=1{,}225$ y $-0{,}4/0{,}4899=-0{,}816$.
  \item Desviación muestral $\sqrt{1{,}2/4}=0{,}5477$; ventajas $1{,}095$ y $-0{,}730$.
  \item $0{,}4^{5}+0{,}6^{5}=0{,}0102+0{,}0778=0{,}0880$.
\end{enumerate}
\resultado \textbf{con la desviación poblacional, $1{,}225$ para las respuestas correctas y $-0{,}816$ para las incorrectas; con la muestral, $1{,}095$ y $-0{,}730$; un $8{,}8\,\%$ de los grupos no aporta señal.}
\comprobacion las ventajas suman cero: $2\cdot 1{,}225-3\cdot 0{,}816=0$.
\repasar{M21, §3.4, §4.1, §4.2, §5.1, §5.3 y §6.3.}
\end{solucion}

% =====================================================================  B.5
\begin{solucion}{eje:p5}
\apartado{(a)}
\datos $N=4$; $\df$: «gato»~2, «negro»~2, el resto~1.
\planteamiento $\idf_t=\ln(N/\df_t)$, pesos $f_{t,d}\idf_t$ y coseno entre los vectores de la consulta y de cada documento (M22, §2.2 y~§2.3).
\pasos
\begin{enumerate}[label=\arabic*.]
  \item $\idf=\ln 2=0{,}6931$ para «gato» y «negro»; $\ln 4=1{,}3863$ para los demás términos.
  \item $q=(0{,}6931;\ 0{,}6931)$ en (gato, negro), con $\norm q=0{,}9803$.
  \item $D_1$: $\norm{D_1}=\sqrt{2\cdot 0{,}4805+1{,}9218}=1{,}6979$; $q\cdot D_1=0{,}9609$; $\cos=0{,}5774$.
  \item $D_2$: «negro» pesa $2\cdot 0{,}6931=1{,}3863$, como «perro» y «corre»; $\norm{D_2}=2{,}4011$; $q\cdot D_2=0{,}9609$; $\cos=0{,}4082$.
  \item $D_3$: $\norm{D_3}=\sqrt{0{,}4805+2\cdot 1{,}9218}=2{,}0794$; $q\cdot D_3=0{,}4805$; $\cos=0{,}2357$.
  \item $D_4$ no comparte términos: $\cos=0$.
\end{enumerate}
\resultado \textbf{$D_1$ ($0{,}577$), $D_2$ ($0{,}408$), $D_3$ ($0{,}236$), $D_4$ ($0$).}
\comprobacion los valores exactos son $1/\sqrt3$, $1/\sqrt6$ y $1/(3\sqrt2)$: el factor $\ln 2$ se cancela.

\apartado{(b)}
\datos longitudes $3$, $4$, $3$ y $2$; $\bar\ell=3$; $k_1=1{,}2$; $b=0{,}75$.
\planteamiento $\BM(q,d)=\sum_t\idf^{\mathrm{BM}}_t\,g(f_{t,d},\ell)$ con $g(f,\ell)=f(k_1+1)/\bigl(f+k_1(1-b+b\ell)\bigr)$ (M22, §2.6).
\pasos
\begin{enumerate}[label=\arabic*.]
  \item Para $n_t=2$: $\idf^{\mathrm{BM}}=\ln(1+2{,}5/2{,}5)=\ln 2=0{,}6931$.
  \item $D_1$ ($\ell=1$): $g(1,1)=1$ para cada término; $\BM=1{,}3863$.
  \item $D_2$ ($\ell=4/3$): $1-b+b\ell=1{,}25$; $g(2,4/3)=4{,}4/(2+1{,}5)=1{,}2571$; $\BM=0{,}8714$.
  \item $D_3$: $\BM=0{,}6931$. $D_4$: $0$.
  \item Con una sola aparición, $g(1,4/3)=2{,}2/2{,}5=0{,}88$ y la puntuación sería $0{,}6100$: la segunda aparición añade $0{,}2614$, menos que la primera ($0{,}6100$). En tf-idf, en cambio, duplica el peso.
\end{enumerate}
\resultado \textbf{$1{,}386$; $0{,}871$; $0{,}693$; $0$: el mismo orden que en (a).}
\comprobacion $g(2,4/3)<k_1+1=2{,}2$, de acuerdo con la saturación.
\criterio 1 punto por las puntuaciones; 1 por el efecto de la saturación y la comparación.

\apartado{(c)}
\datos posiciones relevantes: sistema~1, $\{1,4,9\}$; sistema~2, $\{2,3,6\}$; $\abs{\Rel}=4$.
\planteamiento $\AP=\frac{1}{\abs{\Rel}}\sum_i\Pk{k_i}$; R-precisión $=\Pk{\abs{\Rel}}$; $\RR=1/k_1$ (M22, §3.2).
\pasos
\begin{enumerate}[label=\arabic*.]
  \item Sistema~1: $\Pk{1}=1$, $\Pk{4}=2/4$, $\Pk{9}=3/9$; $\AP=(1+0{,}5+0{,}3333)/4=0{,}4583$.
  \item Sistema~2: $\Pk{2}=1/2$, $\Pk{3}=2/3$, $\Pk{6}=3/6$; $\AP=(0{,}5+0{,}6667+0{,}5)/4=0{,}4167$.
  \item R-precisión ($\Pk{4}$): $2/4=0{,}5$ en ambos. $\Pk{10}=0{,}3$ y $\Rk{10}=0{,}75$ en ambos.
  \item Rango recíproco: $1$ y $1/2$.
\end{enumerate}
\resultado \textbf{AP $0{,}458$ frente a $0{,}417$; R-precisión $0{,}5$, $\Pk{10}=0{,}3$ y $\Rk{10}=0{,}75$ iguales; RR $1$ frente a $0{,}5$.} La AP prefiere el sistema~1 y la R-precisión empata: la AP premia colocar relevantes en las primeras posiciones, y un acierto en la posición~1 pesa mucho.
\comprobacion el cuarto relevante no se recupera y aporta~0 a la AP; por eso ninguna AP supera $3/4$.

\apartado{(d)}
\datos sistema~1; $10\,000$ documentos; el cuarto relevante en una posición entre $11$ y $10\,000$.
\planteamiento la AP crece cuanto antes aparezca el cuarto relevante, que añade $\frac14\cdot\frac{4}{k}$.
\pasos
\begin{enumerate}[label=\arabic*.]
  \item Mejor caso, $k=11$: $\AP=(1{,}8333+4/11)/4=(1{,}8333+0{,}3636)/4=0{,}5492$.
  \item Peor caso, $k=10\,000$: $\AP=(1{,}8333+0{,}0004)/4=0{,}4584$.
  \item El valor de (c), $0{,}4583$, puede diferir del real hasta en $0{,}5492-0{,}4583=0{,}0909=1/11$.
\end{enumerate}
\resultado \textbf{AP entre $0{,}458$ y $0{,}549$; error máximo $0{,}0909$.}
\comprobacion el error máximo coincide con la aportación máxima del relevante que falta, $\frac14\cdot\frac{4}{11}=\frac{1}{11}$.
\repasar{M22, §2.2, §2.3, §2.6 y §3.2.}
\end{solucion}

% =====================================================================  B.6
\begin{solucion}{eje:p6}
\apartado{(a)}
\datos $q=(1,2,2)$; $c_1=(2,1,2)$, $c_2=(0,3,4)$, $c_3=(10,5,0)$.
\planteamiento producto escalar $q\cdot c$ y coseno $q\cdot c/(\norm q\norm c)$ (M22, §4.2).
\pasos
\begin{enumerate}[label=\arabic*.]
  \item $\norm q=3$; $\norm{c_1}=3$, $\norm{c_2}=5$, $\norm{c_3}=\sqrt{125}=11{,}18$.
  \item Productos: $8$, $14$ y $20$. Cosenos: $8/9=0{,}889$, $14/15=0{,}933$ y $20/33{,}54=0{,}596$.
\end{enumerate}
\resultado \textbf{por producto escalar, $c_3$, $c_2$, $c_1$; por coseno, $c_2$, $c_1$, $c_3$.} El producto escalar incluye la norma del fragmento y favorece los vectores largos. Si se normalizan todos los embeddings (y la consulta) a norma~1 antes de indexarlos, el producto escalar coincide con el coseno.
\comprobacion los tres cosenos están en $[-1,1]$ y el de $c_3$ es el menor, aunque su producto escalar sea el mayor.

\apartado{(b)}
\datos $L=5000$, $\tau=512$, $s=64$, $k=4$.
\planteamiento $n=\lceil(L-\tau)/(\tau-s)\rceil+1$; el trozo $j$ empieza en $j(\tau-s)$ (M22, §5.3).
\pasos
\begin{enumerate}[label=\arabic*.]
  \item Paso: $448$. $n=\lceil 4488/448\rceil+1=\lceil 10{,}02\rceil+1=12$.
  \item El último empieza en $11\cdot 448=4928$ y tiene $5000-4928=72$ tokens.
  \item Tokens indexados: $11\cdot 512+72=5704$ (la aproximación $L\tau/(\tau-s)$ da $5714$).
  \item \emph{Prompt}: $4\cdot 512+300=2348$ tokens.
\end{enumerate}
\resultado \textbf{12 fragmentos; el último empieza en la posición 4928; unos 5700 tokens indexados; \emph{prompt} de 2348 tokens.}
\comprobacion el trozo~10 termina en $4480+511=4991<4999$, de modo que el duodécimo es necesario.

\apartado{(c)}
\datos $n=8$; $p=0{,}9$; objetivo $0{,}95$.
\planteamiento éxito $p^{n}$; primer fallo geométrico de media $1/(1-p)$; con reintento, $1-(1-p)^{2}$ por paso (M22, §6.4).
\pasos
\begin{enumerate}[label=\arabic*.]
  \item $0{,}9^{8}=0{,}4305$; primer fallo, en media, en el paso $1/0{,}1=10$.
  \item Con reintento: $1-0{,}1^{2}=0{,}99$ por paso y $0{,}99^{8}=0{,}9227$.
  \item $p^{8}=0{,}95\Rightarrow p=0{,}95^{1/8}=e^{\ln 0{,}95/8}=0{,}99361$.
\end{enumerate}
\resultado \textbf{$0{,}430$; 10 pasos; $0{,}923$ con verificación y reintento; hace falta $p\ge 0{,}9936$.}
\comprobacion $1-8\cdot(1-0{,}99361)=0{,}9489\approx 0{,}95$.

\apartado{(d)}
\datos $\hat p_A=640/800=0{,}8$; $\hat p_B=0{,}76$; $n=800$ cada uno.
\planteamiento $\SE=\sqrt{\hat p(1-\hat p)/n}$; intervalo conservador $\bar x\pm z_{\alpha/2}/(2\sqrt n)$, que al $95\,\%$ es $\bar x\pm 1/\sqrt n$; varianzas que se suman para muestras independientes (M22, §3.3; M23, §6.1).
\pasos
\begin{enumerate}[label=\arabic*.]
  \item $\SE_A=\sqrt{0{,}16/800}=0{,}01414$; $1{,}96\cdot 0{,}01414=0{,}0277$; intervalo $[0{,}772;\ 0{,}828]$.
  \item Conservador: $1/\sqrt{800}=0{,}0354$; intervalo $[0{,}765;\ 0{,}835]$.
  \item $1/\sqrt n=0{,}01\Rightarrow n=10\,000$ (con $1{,}96$ en lugar de~2, $9604$). Al $80\,\%$: $1{,}2816/(2\sqrt n)=0{,}01\Rightarrow n=(64{,}08)^{2}=4106{,}2$, es decir, $4107$ preguntas.
  \item $\SE_{A-B}=\sqrt{0{,}16/800+0{,}1824/800}=\sqrt{4{,}28\cdot 10^{-4}}=0{,}02069$.
  \item Intervalo: $0{,}04\pm 1{,}96\cdot 0{,}02069=0{,}04\pm 0{,}0405=[-0{,}0005;\ 0{,}0805]$; $z=0{,}04/0{,}02069=1{,}93<1{,}96$.
\end{enumerate}
\resultado \textbf{A: $[0{,}772;\ 0{,}828]$ (conservador $[0{,}765;\ 0{,}835]$); $10\,000$ y $4107$ preguntas; la diferencia de 4 puntos no es significativa al $5\,\%$, por muy poco.} Evaluar ambos sistemas sobre las mismas preguntas y analizar la diferencia pregunta a pregunta reduce el error, porque la covarianza positiva entre los aciertos se resta.
\comprobacion el intervalo contiene el~0 por $0{,}0005$, coherente con $z=1{,}93$, apenas por debajo de $1{,}96$.
\criterio 1 punto por (i); 1 por (ii); 2 por (iii), de los que 0,5 son por la propuesta de diseño emparejado.
\repasar{M22, §3.3, §4.2, §5.3 y §6.4; M23, §6.1.}
\end{solucion}

% =====================================================================  B.7
\begin{solucion}{eje:p7}
\apartado{(a)}
\datos $W_1=(1,0)$, $W_2=(-1,0)$, $W_3=(0,1)$, $W_4=(0,-1)$; sin sesgo; probabilidad de activación $0{,}1$.
\planteamiento $(W\T W)_{ij}=W_i\cdot W_j$; $\hat x=\relu(W\T Wx)$; $D_i=\norm{W_i}^{2}/\sum_j(\hat W_i\T W_j)^{2}$ (M23, §3.3 y~§3.4).
\pasos
\begin{enumerate}[label=\arabic*.]
  \item $W\T W=\begin{pmatrix}1&-1&0&0\\-1&1&0&0\\0&0&1&-1\\0&0&-1&1\end{pmatrix}$: los pares antipodales interfieren con $-1$ y los perpendiculares no interfieren.
  \item $x=e_1$: $W\T Wx=(1,-1,0,0)$; $\hat x=(1,0,0,0)$; error~0.
  \item $x=e_1+e_2$: $Wx=0$, luego $\hat x=0$ y el error es $1+1=2$.
  \item $x=e_1+e_3$: $W\T Wx=(1,-1,1,-1)$; $\hat x=(1,0,1,0)$; error~0.
  \item $D_i=1/(1+1+0+0)=1/2$ para cada $i$; $\sum_iD_i=2=m$.
  \item Un par tiene sus dos características activas con probabilidad $0{,}01$; alguno de los dos pares, $1-0{,}99^{2}=0{,}0199$.
\end{enumerate}
\resultado \textbf{errores $0$, $2$ y $0$; $D_i=1/2$ con suma $2$; probabilidad de interferencia $0{,}0199$.} La $\relu$ elimina la interferencia negativa de una característica aislada; el modelo solo falla cuando se activan a la vez las dos de un par antipodal, lo que es raro si las características son dispersas.
\comprobacion la suma de las dimensiones alcanza la cota $\rango W=2$: la configuración aprovecha todo el espacio.
\criterio 1,5 puntos por las reconstrucciones; 1 por $D_i$; 0,5 por la probabilidad.

\apartado{(b)}
\datos $d_1=(1,0)$, $d_2=(0,1)$, $d_3=(0{,}8;\ 0{,}6)$; $x=(1{,}6;\ 1{,}2)=2d_3$; $\lambda=0{,}2$.
\planteamiento pérdida $\norm{x-\hat x}^{2}+\lambda\sum_jf_j\norm{d_j}$; con un solo latente, $\norm{x-f_3d_3}^{2}=(2-f_3)^{2}$ y se aplica el umbral suave (M23, §4.2 y~§4.3).
\pasos
\begin{enumerate}[label=\arabic*.]
  \item $f=(1{,}6;\ 1{,}2;\ 0)$: $\hat x=x$, penalización $0{,}2\cdot 2{,}8=0{,}56$.
  \item $f=(0,0,2)$: $\hat x=2d_3=x$, penalización $0{,}2\cdot 2=0{,}40$.
  \item $\min_{f_3\ge 0}(2-f_3)^{2}+0{,}2f_3$: $f_3^{\ast}=\max(2-0{,}1,\,0)=1{,}9$.
  \item $\hat x=1{,}9\,d_3=(1{,}52;\ 1{,}14)$; error $0{,}08^{2}+0{,}06^{2}=0{,}01$; pérdida $0{,}01+0{,}38=0{,}39$.
\end{enumerate}
\resultado \textbf{$0{,}56$ frente a $0{,}40$: la penalización prefiere el código de un solo latente; el óptimo es $f_3=1{,}9$, con pérdida $0{,}39$.} El sesgo es el \emph{encogimiento}: el autoencoder subestima en $\lambda/2=0{,}1$ la intensidad de la característica activa.
\comprobacion $0{,}39<0{,}40$: aceptar un pequeño error de reconstrucción compensa la reducción de la penalización.

\apartado{(c)}
\datos grupos de 80, 70 y 50 respuestas con confianzas $0{,}95$, $0{,}7$ y $0{,}4$ y aciertos 60, 49 y 25; $t=0{,}75$.
\planteamiento $\ECE=\sum_b\frac{\abs{B_b}}{n}\abs{\mathrm{acc}_b-\mathrm{conf}_b}$; $\MCE=\max_b\abs{\mathrm{acc}_b-\mathrm{conf}_b}$; responder conviene si $p>t$ (M23, §6.3 y~§6.4).
\pasos
\begin{enumerate}[label=\arabic*.]
  \item Exactitudes: $0{,}75$, $0{,}70$ y $0{,}50$; diferencias $0{,}20$, $0$ y $0{,}10$.
  \item $\ECE=0{,}4\cdot 0{,}20+0{,}35\cdot 0+0{,}25\cdot 0{,}10=0{,}105$; $\MCE=0{,}20$.
  \item Confianza media $(76+49+20)/200=0{,}725$; exactitud $134/200=0{,}67$.
  \item Con $t=0{,}75$ el error cuesta $-3$; en el grupo de $0{,}7$, responder vale $0{,}7-0{,}3\cdot 3=-0{,}2<0$.
\end{enumerate}
\resultado \textbf{$\ECE=0{,}105$, $\MCE=0{,}20$; el modelo está sobreconfiado ($0{,}725$ frente a $0{,}67$). En el grupo de $0{,}7$ le conviene abstenerse.}
\comprobacion el criterio general $p>t$ da lo mismo: $0{,}7<0{,}75$.

\apartado{(d)}
\datos $p=0{,}01$, $k=200$; 500 intentos sin éxito.
\planteamiento $1-(1-p)^{k}\ge 1-e^{-kp}$; regla del 3, $p_u\le 3/n$, y cota exacta $1-0{,}05^{1/n}$ (M23, §6.6).
\pasos
\begin{enumerate}[label=\arabic*.]
  \item $1-0{,}99^{200}=0{,}8660$; cota $1-e^{-2}=0{,}8647$.
  \item $3/500=0{,}6\,\%$; exacta $1-0{,}05^{1/500}=0{,}597\,\%$.
\end{enumerate}
\resultado \textbf{$0{,}866$ (cota $0{,}865$); tasa de éxito menor que $0{,}6\,\%$ con confianza del $95\,\%$ (exacta $0{,}597\,\%$).} Un ataque raro en una prueba es casi seguro para un atacante persistente, y cero éxitos no significa tasa nula.
\comprobacion la cota exacta queda por debajo de la de la regla del 3, como garantiza $1-p\le e^{-p}$.
\repasar{M23, §3.3, §3.4, §4.2, §4.3, §6.3, §6.4 y §6.6.}
\end{solucion}
